\documentclass[11pt]{article}

\usepackage[margin=1in]{geometry}

\usepackage[T1]{fontenc}
\usepackage[utf8]{inputenc}
\usepackage{lmodern}
\usepackage{microtype}

% Allow common Unicode diagram glyphs from markdown source under pdfLaTeX.
\DeclareUnicodeCharacter{2502}{|}
\DeclareUnicodeCharacter{25BC}{\ensuremath{\downarrow}}

\usepackage{amsmath,amssymb,amsfonts}
\usepackage{calc}
\usepackage{graphicx}
\usepackage{booktabs}
\usepackage{array}
\usepackage{tabularx}
\usepackage{longtable}
\usepackage{multirow}
\usepackage{caption}
\usepackage{subcaption}
\usepackage{float}
\usepackage{enumitem}
\usepackage{xcolor}
\usepackage{listings}
\usepackage{xurl}
\usepackage[hidelinks]{hyperref}
\usepackage[nameinlink,noabbrev]{cleveref}

\setlength{\parskip}{0.35em}
\setlength{\parindent}{0pt}
\setlength{\emergencystretch}{2em}

\definecolor{listinggray}{gray}{0.97}

\lstset{
  basicstyle=\ttfamily\small,
  backgroundcolor=\color{listinggray},
  columns=fullflexible,
  keepspaces=true,
  breaklines=true,
  breakatwhitespace=false,
  frame=single,
  showstringspaces=false,
  tabsize=2
}

\newcolumntype{Y}{>{\raggedright\arraybackslash}X}
\newcolumntype{C}[1]{>{\centering\arraybackslash}p{#1}}

% Pandoc emits \tightlist for compact markdown lists; define it when absent.
\providecommand{\tightlist}{%
  \setlength{\itemsep}{0pt}\setlength{\parskip}{0pt}}

\title{ProtoGPU: A Semantic GPU ISA Prototyping Framework for Rapid Instruction Exploration}
\author{Han Zhenzhong \\ ProtoGPU Project \\ \texttt{https://github.com/Han-Zhenzhong/ProtoGPU}}



\date{}

\usepackage{color}
\usepackage{fancyvrb}
\newcommand{\VerbBar}{|}
\newcommand{\VERB}{\Verb[commandchars=\\\{\}]}
\DefineVerbatimEnvironment{Highlighting}{Verbatim}{commandchars=\\\{\}}
% Add ',fontsize=\small' for more characters per line
\newenvironment{Shaded}{}{}
\newcommand{\AlertTok}[1]{\textcolor[rgb]{1.00,0.00,0.00}{\textbf{#1}}}
\newcommand{\AnnotationTok}[1]{\textcolor[rgb]{0.38,0.63,0.69}{\textbf{\textit{#1}}}}
\newcommand{\AttributeTok}[1]{\textcolor[rgb]{0.49,0.56,0.16}{#1}}
\newcommand{\BaseNTok}[1]{\textcolor[rgb]{0.25,0.63,0.44}{#1}}
\newcommand{\BuiltInTok}[1]{\textcolor[rgb]{0.00,0.50,0.00}{#1}}
\newcommand{\CharTok}[1]{\textcolor[rgb]{0.25,0.44,0.63}{#1}}
\newcommand{\CommentTok}[1]{\textcolor[rgb]{0.38,0.63,0.69}{\textit{#1}}}
\newcommand{\CommentVarTok}[1]{\textcolor[rgb]{0.38,0.63,0.69}{\textbf{\textit{#1}}}}
\newcommand{\ConstantTok}[1]{\textcolor[rgb]{0.53,0.00,0.00}{#1}}
\newcommand{\ControlFlowTok}[1]{\textcolor[rgb]{0.00,0.44,0.13}{\textbf{#1}}}
\newcommand{\DataTypeTok}[1]{\textcolor[rgb]{0.56,0.13,0.00}{#1}}
\newcommand{\DecValTok}[1]{\textcolor[rgb]{0.25,0.63,0.44}{#1}}
\newcommand{\DocumentationTok}[1]{\textcolor[rgb]{0.73,0.13,0.13}{\textit{#1}}}
\newcommand{\ErrorTok}[1]{\textcolor[rgb]{1.00,0.00,0.00}{\textbf{#1}}}
\newcommand{\ExtensionTok}[1]{#1}
\newcommand{\FloatTok}[1]{\textcolor[rgb]{0.25,0.63,0.44}{#1}}
\newcommand{\FunctionTok}[1]{\textcolor[rgb]{0.02,0.16,0.49}{#1}}
\newcommand{\ImportTok}[1]{\textcolor[rgb]{0.00,0.50,0.00}{\textbf{#1}}}
\newcommand{\InformationTok}[1]{\textcolor[rgb]{0.38,0.63,0.69}{\textbf{\textit{#1}}}}
\newcommand{\KeywordTok}[1]{\textcolor[rgb]{0.00,0.44,0.13}{\textbf{#1}}}
\newcommand{\NormalTok}[1]{#1}
\newcommand{\OperatorTok}[1]{\textcolor[rgb]{0.40,0.40,0.40}{#1}}
\newcommand{\OtherTok}[1]{\textcolor[rgb]{0.00,0.44,0.13}{#1}}
\newcommand{\PreprocessorTok}[1]{\textcolor[rgb]{0.74,0.48,0.00}{#1}}
\newcommand{\RegionMarkerTok}[1]{#1}
\newcommand{\SpecialCharTok}[1]{\textcolor[rgb]{0.25,0.44,0.63}{#1}}
\newcommand{\SpecialStringTok}[1]{\textcolor[rgb]{0.73,0.40,0.53}{#1}}
\newcommand{\StringTok}[1]{\textcolor[rgb]{0.25,0.44,0.63}{#1}}
\newcommand{\VariableTok}[1]{\textcolor[rgb]{0.10,0.09,0.49}{#1}}
\newcommand{\VerbatimStringTok}[1]{\textcolor[rgb]{0.25,0.44,0.63}{#1}}
\newcommand{\WarningTok}[1]{\textcolor[rgb]{0.38,0.63,0.69}{\textbf{\textit{#1}}}}


\begin{document}

\maketitle

\begin{abstract}



The evolution of GPU instruction sets is often constrained by long
hardware design cycles and limited prototyping infrastructure.
Evaluating new GPU instructions typically requires complex
microarchitectural simulators or full hardware implementations, which
significantly slows down the exploration of ISA extensions.

This paper presents \textbf{ProtoGPU}, a semantic GPU execution
framework designed for rapid prototyping and validation of new GPU
instructions. ProtoGPU executes CUDA and PTX programs using a functional
SIMT execution model without requiring cycle-accurate hardware
simulation. The framework introduces an \textbf{IR-based instruction
expansion mechanism} that enables new instructions to be integrated by
mapping them to intermediate representation (IR) operations and
expanding them into micro-operations executed by a semantic SIMT engine.

To demonstrate the capability of the framework, we prototype a
warp-level reduction instruction, \texttt{warp\_reduce\_add}, and
integrate it into the CUDA programming workflow using inline PTX
assembly. The instruction is compiled by Clang into PTX and executed on
ProtoGPU through a CUDA Runtime shim path. Current repository evidence
validates the instruction semantics on focused demos and
unit/integration tests, and the quantitative results reported in this
paper are grounded in reproducible experiment logs from that workflow.

ProtoGPU provides a lightweight infrastructure for GPU ISA research and
hardware--software co-design, enabling rapid exploration of instruction
semantics and compiler integration prior to hardware implementation.

\end{abstract}


\hypertarget{introduction}{%
\section{Introduction}\label{introduction}}

Modern GPUs continue to evolve with increasingly complex instruction
sets to support emerging workloads such as machine learning, graph
analytics, and large-scale data processing. Designing and evaluating new
GPU instructions is challenging due to several factors:

\begin{itemize}
\tightlist
\item
  GPU architectures are complex and largely proprietary.
\item
  Hardware implementation cycles are long and costly.
\item
  Existing simulation frameworks focus primarily on microarchitectural
  performance modeling rather than instruction set exploration.
\end{itemize}

Architecture simulators such as GPGPU-Sim [1] and Accel-Sim [2]
provide detailed modeling of GPU hardware pipelines and memory systems.
However, these simulators are primarily designed for performance
analysis rather than rapid instruction set prototyping. Integrating new
instructions into such frameworks often requires significant
modifications to the simulated hardware pipeline.

For early-stage ISA exploration, researchers often require a lightweight
environment that can:

\begin{enumerate}
\def\labelenumi{\arabic{enumi}.}
\tightlist
\item
  Execute representative GPU programs/demos
\item
  Rapidly integrate new instructions
\item
  Validate instruction semantics
\item
  Evaluate compiler integration
\end{enumerate}

To address these needs, we present \textbf{ProtoGPU}, a semantic GPU
execution framework designed for \textbf{GPU ISA prototyping}.

ProtoGPU executes CUDA and PTX programs using a functional SIMT
execution model. Instead of simulating hardware timing behavior,
ProtoGPU focuses on instruction semantics and program correctness. New
instructions can be rapidly integrated by mapping them to IR operations
that are internally expanded into micro-operations executed by the SIMT
engine.

To demonstrate the usefulness of ProtoGPU, we implement a new warp-level
instruction called \texttt{warp\_reduce\_add}, which performs reduction
across threads within a warp. The instruction is exposed to CUDA
programs through inline PTX assembly and compiled using Clang. ProtoGPU
executes the generated PTX in the current demo workflow and validates
the intended semantics in focused test coverage.

This work makes the following contributions:

\begin{enumerate}
\def\labelenumi{\arabic{enumi}.}
\tightlist
\item
  \textbf{ProtoGPU}, a semantic GPU execution framework capable of
  executing CUDA and PTX programs.
\item
  An \textbf{IR-based instruction prototyping mechanism} that enables
  rapid integration of new GPU instructions.
\item
  A \textbf{CUDA-compatible execution environment} (CUDA Runtime shim +
  PTX override workflow) that allows evaluation of ISA extensions on
  executable CUDA demos.
\item
  A \textbf{case study implementing a warp-level reduction instruction},
  demonstrating the effectiveness of the framework for ISA exploration.
\end{enumerate}


\hypertarget{background}{%
\section{Background}\label{background}}

\hypertarget{gpu-execution-model}{%
\subsection{GPU Execution Model}\label{gpu-execution-model}}

GPUs use a \textbf{SIMT (Single Instruction Multiple Threads)} execution
model in which groups of threads called \textbf{warps} execute the same
instruction simultaneously.

Key concepts include:

\begin{itemize}
\tightlist
\item
  \textbf{Warp}: typically 32 threads executing in lock-step
\item
  \textbf{Thread blocks}: groups of threads scheduled on streaming
  multiprocessors
\item
  \textbf{PTX}: NVIDIA's intermediate GPU assembly language [4]
\item
  \textbf{Warp-level primitives}: operations that allow threads within a
  warp to communicate efficiently
\end{itemize}

Warp-level operations such as shuffle and reduction are commonly used in
GPU kernels.


\hypertarget{challenges-in-gpu-isa-prototyping}{%
\subsection{Challenges in GPU ISA
Prototyping}\label{challenges-in-gpu-isa-prototyping}}

Evaluating new GPU instructions typically requires one of the following
approaches:

\hypertarget{hardware-implementation}{%
\subsubsection{Hardware Implementation}\label{hardware-implementation}}

Direct implementation in hardware provides accurate results but requires
substantial engineering effort and long design cycles.

\hypertarget{microarchitectural-simulation}{%
\subsubsection{Microarchitectural
Simulation}\label{microarchitectural-simulation}}

Cycle-accurate GPU simulators provide detailed performance modeling but
are often complex and difficult to modify.

\hypertarget{dynamic-compilation-frameworks}{%
\subsubsection{Dynamic Compilation
Frameworks}\label{dynamic-compilation-frameworks}}

Some frameworks translate GPU programs to other architectures but do not
directly support instruction set exploration.

ProtoGPU aims to provide a lightweight alternative that focuses on
\textbf{instruction semantics and integration feasibility}.


\hypertarget{protogpu-architecture}{%
\section{ProtoGPU Architecture}\label{protogpu-architecture}}

ProtoGPU implements a semantic GPU execution framework capable of
running CUDA demo programs through a PTX execution pipeline.

\hypertarget{system-overview}{%
\subsection{System Overview}\label{system-overview}}

The execution pipeline of ProtoGPU is shown in Fig.~\ref{fig:protogpu-pipeline}.

\begin{figure}[H]
\centering
\begin{tabular}{c}
\texttt{CUDA Program} \\
$\downarrow$ \\
\texttt{Clang / CUDA Compiler} \\
$\downarrow$ \\
\texttt{PTX Code} \\
$\downarrow$ \\
\texttt{PTX Parser} \\
$\downarrow$ \\
\texttt{Intermediate Representation (IR)} \\
$\downarrow$ \\
\texttt{Micro-Operation Expansion} \\
$\downarrow$ \\
\texttt{SIMT Execution Engine}
\end{tabular}
\caption{ProtoGPU execution pipeline from CUDA source to semantic SIMT execution.}
\label{fig:protogpu-pipeline}
\end{figure}

ProtoGPU operates at the \textbf{semantic execution level}, focusing on
functional correctness rather than microarchitectural timing.


\hypertarget{cuda-integration}{%
\subsection{CUDA Integration}\label{cuda-integration}}

ProtoGPU includes a lightweight \textbf{CUDA compatibility shim} that
allows selected CUDA demo kernels to execute within the framework.

The validated workflow is as follows:

\begin{enumerate}
\def\labelenumi{\arabic{enumi}.}
\tightlist
\item
  CUDA programs are compiled using Clang.
\item
  Inline PTX assembly allows custom instructions to be embedded.
\item
  PTX text is provided to the shim (for example via PTX override) and
  executed by ProtoGPU.
\end{enumerate}

This mechanism allows researchers to evaluate ISA extensions using
executable CUDA demos in a functional (non-cycle-accurate) setting.


\hypertarget{ir-based-instruction-expansion}{%
\subsection{IR-Based Instruction
Expansion}\label{ir-based-instruction-expansion}}

A key feature of ProtoGPU is its \textbf{IR-based instruction
prototyping mechanism}.

Each PTX instruction is mapped to one or more \textbf{IR operations}.
These IR operations are further expanded into \textbf{micro-operations
(uops)} executed by the SIMT engine.

Example mapping is shown in Fig.~\ref{fig:ir-uop-mapping}.

\begin{figure}[H]
\centering
\begin{tabular}{c}
\texttt{PTX Instruction} \\
$\downarrow$ \\
\texttt{IR Operation} \\
$\downarrow$ \\
\texttt{Micro-operations}
\end{tabular}
\caption{Instruction mapping flow from PTX instruction to IR operation and micro-operations.}
\label{fig:ir-uop-mapping}
\end{figure}

This design allows new instructions to be rapidly integrated by defining
IR mappings and micro-operation expansions.


\hypertarget{instruction-prototyping-mechanism}{%
\section{Instruction Prototyping
Mechanism}\label{instruction-prototyping-mechanism}}

Integrating a new instruction into ProtoGPU follows a three-step workflow.

\hypertarget{step-1-define-ptx-isa-mapping}{%
\subsection{Step 1: PTX-ISA
Mapping}\label{step-1-define-ptx-isa-mapping}}

Add the new opcode form to the PTX-ISA mapping asset so frontend
decoding can map it to an IR opcode.

Example:

\begin{lstlisting}
warp_reduce_add
\end{lstlisting}

\hypertarget{step-2-define-descriptor-expansion}{%
\subsection{Step 2: Descriptor
Expansion}\label{step-2-define-descriptor-expansion}}

Add an instruction descriptor entry that expands the IR opcode into one
or more micro-operations.

\hypertarget{step-3-implementenable-micro-operation-semantics}{%
\subsection{Step 3: Micro-Operation
Semantics}\label{step-3-implementenable-micro-operation-semantics}}

Implement (or reuse) execution semantics for the new micro-operation
path in the SIMT execution engine.

This approach allows rapid experimentation with new instruction
semantics.


\hypertarget{case-study-warp_reduce_add}{%
\section{Case Study:
warp\_reduce\_add}\label{case-study-warp_reduce_add}}

To demonstrate the capability of ProtoGPU, we prototype a warp-level
reduction instruction called \texttt{warp\_reduce\_add}.

\hypertarget{motivation}{%
\subsection{Motivation}\label{motivation}}

Warp-level reductions are common in GPU kernels, such as summation and
histogram computations.

Traditional implementations require multiple instructions:

\begin{lstlisting}
shuffle
add
loop
\end{lstlisting}

This sequence introduces additional instruction overhead.


\hypertarget{proposed-instruction}{%
\subsection{Proposed Instruction}\label{proposed-instruction}}

We propose the instruction:

\begin{lstlisting}
warp_reduce_add
\end{lstlisting}

Semantics:

\begin{lstlisting}
r0 = sum of values across threads in the warp
\end{lstlisting}

The instruction performs a warp-level reduction and returns the result
to each participating thread.


\hypertarget{cuda-integration-1}{%
\subsection{CUDA Integration}\label{cuda-integration-1}}

The instruction is exposed to CUDA programs using inline PTX [3,4]:

\begin{Shaded}
\begin{Highlighting}[]
\KeywordTok{asm} \AttributeTok{volatile}\OperatorTok{(}\StringTok{"warp\_reduce\_add.f32 \%0, \%1;"} \OperatorTok{:} \StringTok{"=f"}\OperatorTok{(}\NormalTok{out}\OperatorTok{)} \OperatorTok{:} \StringTok{"f"}\OperatorTok{(}\NormalTok{x}\OperatorTok{));}
\end{Highlighting}
\end{Shaded}

Clang compiles this code and generates PTX; support for custom
instructions in this workflow is validated through ProtoGPU's execution
pipeline and test artifacts [6,10].

ProtoGPU successfully parses and executes the instruction.


\hypertarget{implementation-in-protogpu}{%
\subsection{Implementation in
ProtoGPU}\label{implementation-in-protogpu}}

The instruction integration flow is shown in Fig.~\ref{fig:warp-reduce-add-integration}.

\begin{figure}[H]
\centering
\begin{tabular}{c}
\texttt{warp\_reduce\_add} \\
$\downarrow$ \\
\texttt{IR opcode: warp\_reduce\_add} \\
$\downarrow$ \\
\texttt{EXEC uop: WARP\_REDUCE\_ADD}
\end{tabular}
\caption{Integration path of \texttt{warp\_reduce\_add} from instruction form to IR opcode and execution micro-operation.}
\label{fig:warp-reduce-add-integration}
\end{figure}

The IR operation coordinates data exchange and accumulation across
threads within a warp.


\hypertarget{evaluation}{%
\section{Evaluation}\label{evaluation}}

We evaluate ProtoGPU using focused warp-reduction demos and tests in the
current repository state.

\hypertarget{experimental-setup}{%
\subsection{Experimental Setup}\label{experimental-setup}}

\begin{itemize}
\tightlist
\item
  CUDA kernels compiled using Clang
\item
  PTX executed on ProtoGPU through the CUDA Runtime shim path
\item
  Comparison target: custom \texttt{warp\_reduce\_add} path vs
  shuffle-based reduction PTX path
\end{itemize}

All measurements in Section 6.2 were collected on Linux from the public
ProtoGPU \texttt{v0.1.0} release [10], using the repository shim integration
scripts and shim-exported trace/stats artifacts. The evaluated PTX
targets \texttt{.version\ 6.4} and \texttt{.target\ sm\_70}, and the
custom-op workflow uses a split host-binary/PTX-override path so that
Clang-generated PTX can be executed through the CUDA Runtime shim. These
measurements report functional-semantic execution counters rather than
cycle-level hardware events.


\hypertarget{instruction-count-reduction}{%
\subsection{Instruction Count
Reduction}\label{instruction-count-reduction}}

We report two complementary metrics from shim-side dynamic stats:

\begin{enumerate}
\def\labelenumi{\arabic{enumi}.}
\tightlist
\item
  \textbf{Global dynamic count} (end-to-end kernel-level):
  \texttt{inst.fetch}
\item
  \textbf{Semantics-local equivalent count} (reduction idiom only):
  baseline \texttt{shfl\ +\ add} sequence counters versus
  \texttt{warp\_reduce\_add}
\end{enumerate}

In this experiment, \texttt{inst.commit} matched \texttt{inst.fetch} in
both paths, providing a consistency check for global dynamic instruction
counts.

\hypertarget{global-dynamic-count-end-to-end}{%
\subsubsection{Global dynamic count
(end-to-end)}\label{global-dynamic-count-end-to-end}}

\begin{table}[htbp]
\centering
\caption{Global dynamic count (end-to-end).}
\label{tab:global-dynamic-count}
\begin{tabular}{lccc}
\toprule
Workload & Baseline & Custom & Delta (\%) \\
\midrule
\texttt{warp\_reduce\_add\_demo} kernel & 41 & 24 & -41.5 \\
\bottomrule
\end{tabular}
\end{table}

Baseline is measured using \texttt{inst.fetch}. Evidence artifacts are available in
\texttt{build/test\_out/baseline.stats.json} and
\texttt{build/test\_out/custom.stats.json}.

\hypertarget{semantics-local-equivalent-count-reduction-idiom-only}{%
\subsubsection{Semantics-local equivalent count (reduction idiom
only)}\label{semantics-local-equivalent-count-reduction-idiom-only}}

For the reduction idiom itself, we compare only semantically equivalent
operations:

\begin{itemize}
\tightlist
\item
  Baseline equivalent sequence:
  \texttt{uop.exec.shfl.down} + \texttt{uop.exec.shfl.idx} +
  \texttt{uop.exec.add.f32} = 3 + 1 + 3 = 7
\item
  Custom instruction path: \texttt{uop.exec.warp\_reduce\_add.f32} = 1
\end{itemize}

This yields an \textbf{85.7\% reduction} for the reduction operation
sequence itself.

\begin{table}[htbp]
\centering
\caption{Semantics-local equivalent dynamic count.}
\label{tab:local-dynamic-count}
\begin{tabular}{lccc}
\toprule
Metric scope & Baseline ops & Custom ops & Delta (\%) \\
\midrule
Reduction idiom only & 7 & 1 & -85.7 \\
\bottomrule
\end{tabular}
\end{table}

Artifacts are available in \texttt{build/test\_out/baseline.stats.json} and
\texttt{build/test\_out/custom.stats.json}.

\hypertarget{static-ptx-instruction-line-count-kernel-body}{%
\subsubsection{Static PTX instruction-line count (kernel
body)}\label{static-ptx-instruction-line-count-kernel-body}}

From generated PTX files, we also report static kernel-body
instruction-line counts as a supplementary metric.

\begin{table}[htbp]
\centering
\caption{Static PTX instruction-line count (kernel body).}
\label{tab:static-ptx-count}
\begin{tabular}{lccc}
\toprule
Workload & Baseline & Custom & Delta (\%) \\
\midrule
\texttt{warp\_reduce\_add\_demo} kernel & 40 & 23 & -42.5 \\
\bottomrule
\end{tabular}
\end{table}

Static PTX artifacts are available in
\texttt{build/test\_out/warp\_reduce\_add\_demo\_alternative.ptx} and
\texttt{build/test\_out/warp\_reduce\_add\_demo.ptx}.

Note: the static whole-kernel PTX counting rule counts instruction lines
with explicit opcode + operand formatting and does not include
\texttt{ret;}. Therefore, static totals are one instruction lower than
dynamic \texttt{inst.fetch} totals for both paths; this does not affect
the relative reduction trend.

\hypertarget{static-ptx-semantics-local-equivalent-count}{%
\subsubsection{Static PTX semantics-local equivalent
count}\label{static-ptx-semantics-local-equivalent-count}}

To exclude unrelated kernel instructions, we also compare only the
static PTX instructions implementing the reduction idiom:

\begin{itemize}
\tightlist
\item
  Baseline equivalent PTX sequence:
  \texttt{shfl.sync.*} + \texttt{add.rn.f32} = 4 + 3 = 7
\item
  Custom PTX instruction: \texttt{warp\_reduce\_add.f32} = 1
\end{itemize}

This yields an \textbf{85.7\% static reduction} for the reduction idiom
itself.

\begin{table}[htbp]
\centering
\caption{Semantics-local equivalent static PTX count.}
\label{tab:local-static-ptx-count}
\begin{tabular}{lccc}
\toprule
Metric scope & Baseline ops & Custom ops & Delta (\%) \\
\midrule
Reduction idiom only (static PTX) & 7 & 1 & -85.7 \\
\bottomrule
\end{tabular}
\end{table}

Static PTX artifacts are available in
\texttt{build/test\_out/warp\_reduce\_add\_demo\_alternative.ptx} and
\texttt{build/test\_out/warp\_reduce\_add\_demo.ptx}.

\hypertarget{reproducible-measurement-procedure}{%
\subsubsection{Reproducible measurement
procedure}\label{reproducible-measurement-procedure}}

\begin{enumerate}
\def\labelenumi{\arabic{enumi}.}
\item
  Run the custom instruction path with shim stats export:

\begin{Shaded}
\begin{Highlighting}[]
   \VariableTok{GPUSIM\_TRACE}\OperatorTok{=}\VariableTok{$PWD}\NormalTok{/build/test\_out/custom.trace.jsonl }\DataTypeTok{\textbackslash{}}
   \VariableTok{GPUSIM\_STATS}\OperatorTok{=}\VariableTok{$PWD}\NormalTok{/build/test\_out/custom.stats.json }\DataTypeTok{\textbackslash{}}
   \FunctionTok{bash}\NormalTok{ scripts/run\_cuda\_shim\_e2e\_warp\_reduce\_add\_demo\_cu.sh build}
\end{Highlighting}
\end{Shaded}
\item
  Run the shuffle-based baseline path with shim stats export:

\begin{Shaded}
\begin{Highlighting}[]
   \VariableTok{GPUSIM\_TRACE}\OperatorTok{=}\VariableTok{$PWD}\NormalTok{/build/test\_out/baseline.trace.jsonl }\DataTypeTok{\textbackslash{}}
   \VariableTok{GPUSIM\_STATS}\OperatorTok{=}\VariableTok{$PWD}\NormalTok{/build/test\_out/baseline.stats.json }\DataTypeTok{\textbackslash{}}
   \FunctionTok{bash}\NormalTok{ scripts/run\_cuda\_shim\_e2e\_warp\_reduce\_add\_demo\_alternative\_cu.sh build}
\end{Highlighting}
\end{Shaded}
\item
  Read \texttt{inst.fetch} and \texttt{inst.commit} from the generated
  stats JSON files for the global end-to-end metric.
\item
  Read \texttt{uop.exec.shfl.down}, \texttt{uop.exec.shfl.idx},
  \texttt{uop.exec.add.f32}, and \texttt{uop.exec.warp\_reduce\_add.f32}
  for the semantics-local equivalent metric.
\item
  Report both metrics: global dynamic count and semantics-local
  equivalent count.
\item
  For static PTX kernel-body counting, count instruction lines inside
  the kernel entry body in
  \url{build/test_out/warp_reduce_add_demo_alternative.ptx} and
  \url{build/test_out/warp_reduce_add_demo.ptx}.
\item
  For static PTX semantics-local counting, count only
  \texttt{shfl.sync.*} and \texttt{add.rn.f32} in the baseline PTX
  kernel body, and \texttt{warp\_reduce\_add.f32} in the custom PTX
  kernel body.
\end{enumerate}


\hypertarget{program-simplification}{%
\subsection{Program Simplification}\label{program-simplification}}

Using \texttt{warp\_reduce\_add} simplifies kernel implementations in
focused warp-reduction demos. Quantitative synchronization/counter
reductions should be reported only with attached experiment artifacts.


\hypertarget{limitations}{%
\subsection{Limitations}\label{limitations}}

The evaluation in this paper is intentionally focused on a
warp-reduction bring-up path rather than a broad workload suite. As a
result, the reported instruction-count reductions should be interpreted
as evidence that ProtoGPU can support semantically correct ISA
prototyping and focused quantitative comparison, not as a claim of
general performance improvement across arbitrary CUDA workloads. Broader
external validation would require additional kernels, more instruction
classes, and larger benchmark coverage.


\hypertarget{related-work}{%
\section{Related Work}\label{related-work}}

Existing GPU simulation frameworks primarily focus on microarchitectural
modeling [1,2,8,9].

\hypertarget{gpu-architecture-simulators}{%
\subsection{GPU Architecture
Simulators}\label{gpu-architecture-simulators}}

Cycle-accurate GPU simulators provide detailed performance analysis but
are complex to modify for ISA prototyping [1,2,8,9].

\hypertarget{dynamic-ptx-frameworks}{%
\subsection{Dynamic PTX Frameworks}\label{dynamic-ptx-frameworks}}

Dynamic compilation frameworks translate PTX programs to alternative
execution environments [7].

ProtoGPU differs from these systems by focusing on \textbf{rapid
instruction integration and semantic validation}.


\hypertarget{conclusion}{%
\section{Conclusion}\label{conclusion}}

This paper presented \textbf{ProtoGPU}, a semantic GPU execution
framework for rapid prototyping of GPU ISA extensions.

ProtoGPU enables researchers to integrate new instructions, execute CUDA
demos through the shim workflow, and validate instruction semantics
without requiring cycle-accurate hardware simulation.

Through a focused case study implementing the \texttt{warp\_reduce\_add}
instruction, we show that ProtoGPU can support rapid ISA bring-up,
compiler-adjacent integration, and semantic validation in a functional
SIMT execution setting.

Future work includes extending the framework to support additional ISA
extensions, broadening workload coverage, and exploring compiler-driven
instruction generation.


\hypertarget{references}{%
\section*{References}\label{references}}

\begin{enumerate}[label={[\arabic*]}]
\item
  A. Bakhoda, G. L. Yuan, W. W. L. Fung, H. Wong, and T. M. Aamodt,
  ``Analyzing CUDA Workloads Using a Detailed GPU Simulator,'' in
  \emph{2009 IEEE International Symposium on Performance Analysis of
  Systems and Software (ISPASS)}, 2009.
\item
  M. Khairy, J. Shen, T. M. Aamodt, and T. G. Rogers, ``Accel-Sim: An
  Extensible Simulation Framework for Validated GPU Modeling,'' in
  \emph{Proceedings of the 47th Annual International Symposium on
  Computer Architecture (ISCA)}, 2020.
\item
  NVIDIA Corporation, ``CUDA C++ Programming Guide (CUDA 10.1),''
  \url{https://docs.nvidia.com/cuda/archive/10.1/cuda-c-programming-guide/},
  accessed 2026-04-18.
\item
  NVIDIA Corporation, ``Parallel Thread Execution ISA Version 6.4,''
  \url{https://docs.nvidia.com/cuda/archive/10.1/pdf/ptx_isa_6.4.pdf},
  accessed 2026-04-18.
\item
  C. Lattner and V. Adve, ``LLVM: A Compilation Framework for Lifelong
  Program Analysis and Transformation,'' in \emph{Proceedings of the
  International Symposium on Code Generation and Optimization (CGO)},
  2004.
\item
  LLVM Project, ``Compiling CUDA with Clang (LLVM CUDA Support),''
  \url{https://llvm.org/docs/CompileCudaWithLLVM.html}, accessed 2026-04-18.
\item
  G. Diamos, A. Kerr, S. Yalamanchili, and N. Clark, ``Ocelot: A Dynamic
  Optimization Framework for Bulk-Synchronous Applications in
  Heterogeneous Systems,'' in \emph{Proceedings of the 19th
  International Conference on Parallel Architectures and Compilation
  Techniques (PACT)}, 2010.
\item
  J. Power, J. Hestness, M. S. Orr, M. D. Hill, and D. A. Wood,
  ``gem5-gpu: A Heterogeneous CPU-GPU Simulator,'' \emph{IEEE Computer
  Architecture Letters}, vol.~13, no. 1, pp.~34--36, 2014.
\item
  R. Ubal, J. Sahuquillo, S. Petit, and P. Lopez, ``Multi2Sim: A
  Simulation Framework to Evaluate Multicore-Multithreaded Processors,''
  in \emph{Proceedings of the 16th International Symposium on Computer
  Architecture and High Performance Computing (SBAC-PAD)}, 2007.
\item
  H. Zhenzhong, ``ProtoGPU v0.1.0,'' GitHub, Apr.~2026. [Online].
  Available:
  \url{https://github.com/Han-Zhenzhong/ProtoGPU/releases/tag/v0.1.0}
\end{enumerate}


\end{document}